\documentclass[a4paper,12pt]{article}
\usepackage[utf8]{inputenc}
\usepackage[T1]{fontenc}
\usepackage[heading=true]{ctex} % 支持中文
\usepackage{geometry}
\usepackage{graphicx}
\usepackage{hyperref}
\usepackage{listings}
\usepackage{xcolor}
\usepackage{enumitem}
\pagestyle{plain}

% 页面设置
\geometry{left=2.5cm, right=2.5cm, top=2.5cm, bottom=2.5cm}

% 代码块样式设置
\lstset{
    basicstyle=\small\ttfamily,
    keywordstyle=\color{blue},
    commentstyle=\color{gray},
    stringstyle=\color{red},
    breaklines=true,
    frame=single,
    numbers=left,
    numberstyle=\tiny\color{gray}
}

% 标题信息
\title{\textbf{基于深度学习的多视角视频拼接系统说明文档}}
\author{姓名：张宗桪 \quad 学号：2023K8009991013} % 请在此修改姓名学号
\date{\today}

\begin{document}

\maketitle

\section{原论文原理与实现思路}

本系统基于 Nie Lang 等人提出的深度图像拼接框架 UDIS2，对应的核心论文为 \textit{Parallax-Tolerant Unsupervised Deep Image Stitching (ICCV 2023)}。

\subsection{核视差问题 (Parallax)}
传统的图像拼接方法（如基于特征点匹配和单应性矩阵 Homography 的方法）通常假设场景位于同一平面或相机仅进行纯旋转运动。然而，在监控视频等实际场景中，摄像机往往存在平移，且场景具有深度差异，导致视差现象。传统方法在处理视差时容易产生重影或几何错位。

\subsection{UDIS2 的实现思路}
该论文提出了一种无监督的粗到细拼接框架，主要包含两个核心模块：

\begin{enumerate}
    \item \textbf{变形网络 }：
    不同于计算全局单应性矩阵，该网络预测一个密集的变形场。它通过将图像划分成网格，对每个网格进行自适应的变形，从而灵活地对齐具有不同深度的前景和背景。这种方法具有很强的视差容忍度。
    
    \item \textbf{合成网络 }：
    在变形对齐后，网络并不简单地进行像素平均或线性融合，而是通过学习一个边界掩码，智能地选择保留哪张图像的像素，从而消除拼接缝隙和潜在的伪影。
\end{enumerate}

整个训练过程是无监督的，无需人工标注的拼接真值，这使得该方法具有极强的泛化能力。

\section{视频拼接流程设计}

本项目在 UDIS2 图像拼接模型的基础上，构建了一套完整的视频流处理管线。由于显存限制及视频时序特性，系统采用了"拆帧-级联拼接-重组"的策略，具体流程如下：

\subsection{数据预处理}
首先，通过开发的 \texttt{video\_to\_frames.py} 脚本，将输入的三个监控视频（Video 4, 3, 2）分解为独立的帧序列。
\begin{itemize}
    \item 为了适配 UDIS2 模型的输入要求并防止显存溢出，所有帧被统一调整为 $512 \times 512$ 分辨率。
    \item 视频帧被保存至 \texttt{data/video\_frames/} 目录下的独立子文件夹中 。
\end{itemize}

\subsection{第一阶段：左侧与中间视角拼接 (Input 4 + Input 3)}
通过 \texttt{test\_output\_batch.py} 对 Video 4（左）和 Video 3（中）进行批量推理 。
\begin{enumerate}
    \item \textbf{Warp阶段}：计算两组视频帧之间的变形场。为了提高效率，采用了 \texttt{batch\_size=100} 的批处理设置。
    \item \textbf{Composition阶段}：利用生成的变形特征进行深度融合，生成第一阶段的中间全景视频帧（Video 4-3）。
\end{enumerate}

\subsection{第二阶段：级联拼接 (Result 4-3 + Input 2)}
为了实现三视频全景，本项目采用了级联策略 ：
\begin{enumerate}
    \item \textbf{数据准备}：通过 \texttt{prepare\_second\_stitch.py} 脚本，将第一阶段的拼接结果作为新的"左视图"，将原始的 Video 2（右）作为"右视图"，构建新的输入数据集 \texttt{prepared\_432}。
    \item \textbf{二次推理}：重复 Warp 和 Composition 过程，将 Video 2 融合进之前的全景视野中，生成最终的三视角全景帧。
\end{enumerate}

\subsection{视频重组}
最后，使用 \texttt{create\_video\_from\_frames.py} 脚本将最终生成的图片序列按照 30FPS 的帧率重新编码为 .mp4 视频文件，完成从多视角监控到全景视频的转换 。

\section{代码修改与功能实现}

为了将针对静态图像的 UDIS2 代码库应用于长视频拼接任务，我对源代码进行了以下关键修改：

\begin{enumerate}
    \item \textbf{批处理支持}：在 \texttt{test\_output\_batch.py} 中，添加了对批量图像输入的支持，允许同时处理多帧以提升推理效率。
    \item \textbf{视频拆帧与重组脚本}：新增 \texttt{video\_to\_frames.py} 和 \texttt{create\_video\_from\_frames.py} 脚本，实现视频与帧序列之间的转换。
    \item \textbf{级联拼接逻辑}：编写 \texttt{prepare\_second\_stitch.py} 脚本，自动化处理两阶段拼接的数据准备工作。
    \item \textbf{推理优化}：新增 \texttt{demo.py} 脚本，以适应当前任务需求输入。
\end{enumerate}

\section{系统局限性与改进方向}

尽管系统成功实现了视频拼接，但仍存在以下局限性：

% 【请在此处填写缺陷，例如：】
\begin{itemize}
    \item \textbf{显存限制}：当前处理时受限于显存容量，未来可考虑引入分块处理或多GPU并行计算以提升性能。
    \item \textbf{拼接质量}：在某些复杂场景下，仍出现轻微的重影或拼接缝隙。
    \item \textbf{实时性}：当前实现主要面向离线处理。
\end{itemize}

\section{参考文献与资源}

\begin{enumerate}[label={[\arabic*]}]
    \item Nie, Lang, et al. "Parallax-Tolerant Unsupervised Deep Image Stitching." \textit{Proceedings of the IEEE/CVF International Conference on Computer Vision (ICCV)}. 2023.
    \item GitHub Repository: UDIS2 - \url{https://github.com/nie-lang/UDIS2}
\end{enumerate}

\end{document}